\documentclass[10pt,a4paper]{amsart}
\usepackage[margin=19mm]{geometry}
\usepackage{amsmath,amssymb,mathtools}
\usepackage[scr=rsfs,cal=euler]{mathalfa}
\usepackage{xfrac}
\usepackage[unicode,hidelinks,bookmarks=false]{hyperref}
\newcommand{\ie}{i.e.\ }
\newcommand{\cf}{cf.\ }
\newcommand{\resp}{respectively\ }
\newcommand{\Subsection}{\subsection}
\providecommand{\nobreakdash}{\nobreak\hbox{-}}
\setlength{\emergencystretch}{2em}
\multlinegap0pt
\usepackage{bm}
\usepackage{bbm}
\usepackage{dsfont}
\usepackage{xspace}
\usepackage{cancel}
\usepackage{mathtools}
\usepackage{amsthm}
\makeatletter
\@ifundefined{assumption}{\newtheorem{assumption}{Assumption}}{}
\@ifundefined{lemma}{\newtheorem{lemma}{Lemma}}{}
\@ifundefined{proposition}{\newtheorem{proposition}{Proposition}}{}
\makeatother
\title[On empirical Hodge Laplacians under the manifold hypothesis]{On empirical Hodge Laplacians under the manifold hypothesis}
\author{Jan-Paul Lerch, Martin Wahl, Petr Zamolodtchikov}
\date{Source: arXiv:2504.03427v2 (CC BY 4.0)}
\begin{document}
\maketitle

\noindent\emph{Source and licence.} On empirical Hodge Laplacians under the manifold hypothesis. Version arXiv:2504.03427v2 (CC BY 4.0), distributed under the Creative Commons Attribution 4.0 International licence, as recorded in the arXiv licence metadata for this exact version and shown on the abstract page https://arxiv.org/abs/2504.03427v2. Full rights notice: licenses/arxiv-2504.03427-CC-BY-4.0.txt.\par\smallskip

\noindent\textit{Editorial scope.} Counted unit: the Lemma whose source label is \texttt{lemma:bounds:apply:U:concentation} in the article (source0.tex lines 666--676, source label \texttt{lemma:bounds:apply:U:concentation}), delivered together with the article's own complete original proof of it (source0.tex lines 678--719, one \texttt{proof} environment of 42 lines). No mathematics is written, repaired or re-derived: statement and proof are the article's own text line for line. Required context retained uncounted: Ass:MH (lines 425--426); Ass:SP (lines 446--451); Prop:conc:U:statistic (lines 601--608); lemma:bound:kernel:smoothness (lines 641--651). Every definition, display and label of those lines is retained unchanged. Omitted from this package and not redistributed: 1--424, 427--445, 452--600, 609--640, 652--665, 677--677, 720--944. Nothing inside a retained excerpt is omitted or altered: every retained body is delivered exactly as the article's lines are written, so no omission receipt of any kind accompanies this package. Citations: the retained excerpts carry no citation command at all, so no bibliography entry of the article is transcribed and no reference is redistributed. Reference adaptation: none was needed and none was made; every reference inside the delivered unit resolves inside it, so no unresolved reference or citation is produced. Printed slips kept literal: no printed word, symbol, hypothesis or number of the counted statement or of its proof was altered. Layout and class: the article's own class is \texttt{article} with packages including amsmath, amssymb, amsthm, amsfonts, amstext, amsbsy, amscd, babel, inputenc, csquotes, mathrsfs, bbm, dsfont, xcolor; the wrapper is the lane's pinned \texttt{amsart} editorial wrapper at 10pt on A4, which restates the theorem-like environments the retained text needs and adds the article's own macro definitions that the retained text needs (none), with extra wrapper packages (bm, bbm, dsfont, xspace, cancel, mathtools). The delivered numbering is the unit's own. Assets: no retained span contains a figure environment, a graphics-inclusion command, a PDF-page inclusion, a table or a TikZ picture; no image file is copied into this package, the package declares no asset dependency and no image file is redistributed with it. Licence: version arXiv:2504.03427v2 of this article is distributed under the Creative Commons Attribution 4.0 International licence, as recorded in the arXiv licence metadata for this exact version and shown on the abstract page https://arxiv.org/abs/2504.03427v2. A full rights notice is delivered at licenses/arxiv-2504.03427-CC-BY-4.0.txt. Characters: every retained excerpt is pure ASCII, so the delivered engine, XeLaTeX with its default Latin Modern text font, reports no missing character.
\begin{assumption}[Manifold hypothesis]\label{Ass:MH} Let $X_1,\dots,X_n$ be independent and identical distributed random variables uniformly distributed in a closed and connected submanifold $\mathcal{M}\subseteq\mathbb{R}^p$ with $\dim(\mathcal{M})=d$ and $\operatorname{vol}(\mathcal{M})=1$.
\end{assumption}

\begin{assumption}[Smoothness properties]\label{Ass:SP} There is a constant $C_2>0$ such that
    \begin{align*}
        \|f_a\|_{L^\infty},\|\Delta (f_af_b)\|_{L^\infty},\Big\|\Big(\frac{e^{-t\Delta}-I+t\Delta}{t^2}\Big)(f_af_b)\Big\|_{L^\infty}\leq C_2,
    \end{align*}
    for all $0\leq a,b\leq \ell$, where $f_0\equiv 1$.
\end{assumption}

\begin{proposition}\label{Prop:conc:U:statistic}
    Let $t\in(0,1]$, $A>1$, and $f_1,\dots,f_\ell\in C^\infty(\mathcal{M})$. Suppose that $n\geq 2+2\ell$ and that Assumptions \ref{Ass:MH}--\ref{Ass:SP} are satisfied. Then, with probability at least $1-n^{-A}$, we have
    \begin{align*}
        \Big| U_{n}(\ell,t)&-\frac{1}{t^\ell}\int_{\mathcal{M}^{\ell+1}}(D(f_1,\dots,f_\ell,x_0,\dots,x_\ell))^2\Big(\prod_{b=1}^\ell k_t(x_0,x_b)dx_b\Big)dx_0 \Big|\\
        &\leq C\sum_{j=1}^{\ell+1}\Big(\frac{(\log n)^{j/2}}{t^{d(j-1)/4}n^{j/2}}+\frac{(\log n)^{(j+1)/2}}{t^{d(j-1)/2}n^{(j+1)/2}}\Big),
    \end{align*}
    where $C$ is a constant depending only on $\ell, A, c_1,C_1$ and $C_2$.
\end{proposition}

\begin{lemma}\label{lemma:bound:kernel:smoothness}
    Under the assumptions of Proposition \ref{Prop:conc:U:statistic}, we have 
    \begin{align*}
        \frac{1}{t}k_t(x,y)(f_b(x)-f_b(y))^2\leq C\frac{1}{t^{d/2}}
    \end{align*}
    and
    \begin{align*}
        \int_{\mathcal{M}}\frac{1}{t}k_t(x,y)(f_b(x)-f_b(y))^2\,dy\leq C
    \end{align*}
    for all $x,y\in\mathcal{M}$, all $b=1,\dots,\ell$, and all $t\in(0,1]$, where $C$ is a constant depending only on $c_1,C_1,C_2$.
\end{lemma}

\begin{lemma}\label{lemma:bounds:apply:U:concentation}
    Suppose that the assumptions of Proposition \ref{Prop:conc:U:statistic} are satisfied. Let $J\subseteq \{0,\dots,\ell\}$ be a nonempty subset. Then we have
    \begin{align}\label{eq:bounds:apply:U:concentation:1}
    \int_{\mathcal{M}^{|J^\complement|}}\Big(\frac{1}{t^\ell}\prod_{b=1}^\ell k_t(x_0,x_b)\Big) D^2(f_1,\dots,f_\ell,x_0,\dots,x_\ell)\, dx_{J^\complement}\leq C\frac{1}{t^{d(|J|-1)/2}}    
    \end{align}
    for all $(x_b)_{b\in J}$ and
    \begin{align}\label{eq:bounds:apply:U:concentation:2}
    &\Big(\int_{\mathcal{M}^{|J|}}\Big(\int_{\mathcal{M}^{|J^\complement|}}\Big(\frac{1}{t^\ell}\prod_{b=1}^\ell k_t(x_0,x_b)\Big) D^2(f_1,\dots,f_\ell,x_0,\dots,x_\ell)\, dx_{J^\complement}\Big)^2dx_J\Big)^{1/2}\leq C\frac{1}{t^{d(|J|-1)/4}},
    \end{align}
    where $C>0$ is a constant depending only on $c_1,C_1,C_2$ and $\ell$. Here, $dx_{J^\complement}$ means $\prod_{b\in J^\complement}dx_b$ and $dx_J$ means $\prod_{b\in J}dx_b$.
\end{lemma}

\begin{proof}
    Using the Leibniz formula applied to the transpose, we have
    \begin{align}\label{eq:proof:bounds:apply:U:concentation:1}
        \begin{split}
        D^2(f_1,\dots,f_\ell,x_0,\dots,x_\ell) &=\sum_{\sigma,\tau\in S_\ell}\prod_{b=1}^\ell (f_{\sigma(b)}(x_b)-f_{\sigma(b)}(x_0))(f_{\tau(b)}(x_b)-f_{\tau(b)}(x_0))\\
        &\leq \ell! \sum_{\sigma\in S_\ell}\prod_{b=1}^\ell (f_{\sigma(b)}(x_b)-f_{\sigma(b)}(x_0))^2,
        \end{split}
    \end{align}
    where we used the inequality $xy\leq (x^2+y^2)/2$ and the symmetry in $\sigma,\tau\in S_\ell$. We now consider separately the two cases $0\in J$ and $0\notin J$. 
    
    First, let $0\in J$. Inserting \eqref{eq:proof:bounds:apply:U:concentation:1} into \eqref{eq:bounds:apply:U:concentation:1} and using the Fubini theorem, we get
    \begin{align}
        &\int_{\mathcal{M}^{|J^\complement|}}\Big(\frac{1}{t^\ell}\prod_{b=1}^\ell k_t(x_0,x_b)\Big)\cdot D^2(f_1,\dots,f_\ell,x_0,\dots,x_\ell)\, dx_{J^\complement}\label{eq:proof:bounds:apply:U:concentation:2}\\
        &\leq \ell !\sum_{\sigma\in S_\ell}\int_{\mathcal{M}^{|J^\complement|}}\prod_{b=1}^\ell\frac{1}{t} k_t(x_0,x_b) (f_{\sigma(b)}(x_b)-f_{\sigma(b)}(x_0))^2 dx_{J^\complement}\nonumber\\
        &=\ell !\sum_{\sigma\in S_\ell} \prod_{\substack{b\in J\\ b\neq 0}}\frac{1}{t} k_t(x_0,x_b) (f_{\sigma(b)}(x_b)-f_{\sigma(b)}(x_0))^2\cdot\prod_{b\in J^\complement}\int_{\mathcal{M}}\frac{1}{t} k_t(x_0,x_b) (f_{\sigma(b)}(x_b)-f_{\sigma(b)}(x_0))^2dx_b\Big).\nonumber
    \end{align}
    Inserting Lemma \ref{lemma:bound:kernel:smoothness} the first claim follows in this case. Similarly, 
    \begin{align*}
        &\int_{\mathcal{M}^{|J|}}\Big(\int_{\mathcal{M}^{|J^\complement|}}\Big(\frac{1}{t^\ell}\prod_{b=1}^\ell k_t(x_0,x_b)\Big)\cdot D^2(f_1,\dots,f_\ell,x_0,\dots,x_\ell)\, dx_{J^\complement}\Big)^2dx_J\\
        & \leq (\ell!)^3 \sum_{\sigma\in S_\ell}\int_{\mathcal{M}^{|J|}}\Big(\int_{\mathcal{M}^{|J^\complement|}} \prod_{b=1}^\ell\frac{1}{t} k_t(x_0,x_b) (f_{\sigma(b)}(x_b)-f_{\sigma(b)}(x_0))^2\, dx_{J^\complement}\Big)^2dx_J.
    \end{align*}
    Proceeding as in \eqref{eq:proof:bounds:apply:U:concentation:2}, applying Lemma \ref{lemma:bound:kernel:smoothness} and using the fact that we can also integrate with respect to $x_J$,  the second claim follows. 
    
    Second, let $0\in J^\complement$. Moreover, let $a\in J$ be arbitrary but fixed. Then 
    \begin{align*}
        \int_{\mathcal{M}^{|J^\complement|}}
         \Big( \frac{1}{t^\ell}\prod_{b=1}^\ell k_t(x_0,x_b)\Big)&
        \cdot D^2(f_1,\dots,f_\ell,x_0,\dots,x_\ell)\, dx_{J^\complement}
        \\\leq \ell !\sum_{\sigma\in S_\ell}\int_{\mathcal{M}}\Big[&\prod_{\substack{b\in J\setminus\{a\}}}\frac{1}{t} k_t(x_0,x_b) (f_{\sigma(b)}(x_b)-f_{\sigma(b)}(x_0))^2  
        \\\cdot&\prod_{b\in J^c\setminus\{0\}}\int_{\mathcal{M}}\frac{1}{t} k_t(x_0,x_b) (f_{\sigma(b)}(x_b)-f_{\sigma(b)}(x_0))^2dx_b\\
          \cdot & \frac{1}{t} k_t(x_0,x_a) (f_{\sigma(a)}(x_a)-f_{\sigma(a)}(x_0))^2\Big]\,dx_0
    \end{align*}
    Inserting Lemma \ref{lemma:bound:kernel:smoothness}, the first claim follows in this case, ensuring that we integrate with respect to the $x_0$ in the last step. Similarly,
    \begin{align*}
        &\int_{\mathcal{M}^{|J|}}  \Big[ \int_{\mathcal{M}^{|J^\complement|}} \Big( \frac{1}{t^\ell}\prod_{b=1}^\ell k_t(x_0,x_b) \Big) \cdot D^2(f_1,\dots,f_\ell,x_0,\dots,x_\ell)\, dx_{J^\complement} \Big]^2dx_J\\
        & \leq (\ell!)^3 \sum_{\sigma\in S_\ell}\int_{\mathcal{M}^{|J|}} \Big[\int_{\mathcal{M}^{|J^\complement|}} \Big( \prod_{b=1}^\ell\frac{1}{t} k_t(x_0,x_b) (f_{\sigma(b)}(x_b)-f_{\sigma(b)}(x_0))^2\Big) \, dx_{J^\complement}\Big]^2dx_J\\
        & \leq C(\ell!)^3 \sum_{\sigma\in S_\ell} \int_{\mathcal{M}^{|J|+1}}
        \Big[\int_{\mathcal{M}^{|J^\complement|-1}} \prod_{\substack{b=1\\b\neq a}}^\ell\frac{1}{t} k_t(x_0,x_b) (f_{\sigma(b)}(x_b)-f_{\sigma(b)}(x_0))^2\, dx_{J^\complement\setminus\{0\}}\Big]^2\\
        &\qquad \qquad\qquad\qquad\quad   \cdot\frac{1}{t} k_t(x_0,x_a) (f_{\sigma(a)}(x_a)-f_{\sigma(a)}(x_0))^2dx_0dx_J,
    \end{align*} 
where we applied the Cauchy-Schwarz inequality and Lemma \ref{lemma:bound:kernel:smoothness} in the last inequality. Now, we can proceed as in the first case.
\end{proof}

\end{document}
